\chapter{Compact Groups}

\begin{center}\itshape
On the circle, the functions $1,z,z^{-1},z^2,z^{-2},\ldots$ play two roles. They are one-dimensional representations, and they are the building blocks of Fourier series. What replaces them on the unitary group? The answer will let us integrate a matrix trace by counting permutations.
\end{center}

For a finite group, averaging means adding a function over the group and dividing by its order. For a circle, it means integrating once around the circumference. Haar measure lets us do the same thing on a compact group. Once that average is available, the arguments about finite-group representations from Volume~3 have room to work again.

Issai Schur used invariant integration to study representations in 1924. Hermann Weyl's papers of 1925 and 1926 developed the character theory of compact semisimple Lie groups. Weyl and Fritz Peter then proved their completeness theorem in 1927. Orthogonality was only half the problem: one also had to show that the matrix entries of representations could approximate an arbitrary function.

Weyl later wrote: ``In these days the angel of topology and the devil of abstract algebra fight for the soul of each individual mathematical domain.'' Here the two subjects help each other. Topology supplies a finite invariant measure. Algebra supplies the functions against which we integrate. By the end of the chapter, that integral will be a count of Young tableaux.

Throughout the chapter, $G$ is a compact matrix Lie group: a closed, compact subgroup of $\mathrm{GL}_N(\C)$. The examples to keep in mind are $U(n)$, $SU(n)$ and $SO(n)$. Its Haar probability measure $\mu$ is invariant under multiplication on either side; Volume~9, Theorem~8.3 and Corollary~8.2 establish this. A representation means a continuous homomorphism $\rho\colon G\to\mathrm{GL}(V)$ on a finite-dimensional complex vector space. The zero representation is allowed; an irreducible representation is nonzero. When $U(n)$ appears, $n$ is a positive integer.

\section{Representations of compact groups}

A rotation of the plane preserves the usual inner product. A general matrix representation need not preserve the inner product we first choose. Averaging that inner product removes the dependence on the choice.

\begin{proposition}\label{prop:unitarytrick}
For every representation of $G$, there is an inner product on its space for which the representation is unitary. Every such representation is a direct sum of irreducible representations.
\end{proposition}

\begin{proof}
Choose an inner product $(\cdot,\cdot)_0$ on $V$. Define a new one by
\[
(v,w)_G=\int_G(\rho(g)v,\rho(g)w)_0\,\dd\mu(g).
\]
The integrand is continuous, so each integral exists. Sesquilinearity and conjugate symmetry pass through the integral. If $v\ne0$, then $\rho(g)v\ne0$ for every $g$. The continuous function $g\mapsto\|\rho(g)v\|_0^2$ has a positive minimum on the compact set $G$. Its integral is therefore positive. This proves positive definiteness.

For $h\in G$, substitute $g'=gh$ in the integral. Right invariance gives
\[
(\rho(h)v,\rho(h)w)_G
=\int_G(\rho(gh)v,\rho(gh)w)_0\,\dd\mu(g)
=(v,w)_G.
\]
Thus each $\rho(h)$ is unitary for this inner product.

If $W$ is invariant, then $W^\perp$ is invariant too: for $v\in W^\perp$ and $w\in W$,
$(\rho(h)v,w)_G=(v,\rho(h^{-1})w)_G=0$.
Choose an irreducible subspace of a nonzero $V$ by taking an invariant subspace of least positive dimension. Split off its orthogonal complement and continue. The dimension decreases at every step. If $V=0$, the empty direct sum already gives the assertion.
\end{proof}

Fix orthonormal bases for these invariant inner products. Then $\rho(g^{-1})=\rho(g)^*$. The functions $g\mapsto\rho_{ij}(g)$ are called \emph{matrix coefficients}. They are continuous, and thus belong to $L^2(G,\mu)$.

Schur's lemma from Volume~3 applies without a finiteness assumption on the group: its proof concerns invariant subspaces and eigenvalues of a linear map on a finite-dimensional complex space. Averaging a linear map will turn it into just such an intertwining map.

\begin{theorem}[Schur orthogonality]\label{thm:schurorth}
Choose one unitary representative of each isomorphism class of irreducible representations. If $\rho,\sigma$ are two of these representatives, then
\[
\int_G\rho_{ij}(g)\overline{\sigma_{kl}(g)}\,\dd\mu(g)
=\begin{cases}
 d_\rho^{-1}\delta_{ik}\delta_{jl},&\rho=\sigma,\\
 0,&\rho\not\cong\sigma,
\end{cases}
\]
where $d_\rho=\dim V_\rho$. In particular, the irreducible characters $\chi_\rho=\operatorname{tr}\rho$ are orthonormal in $L^2(G,\mu)$. Two finite-dimensional representations with the same character are isomorphic.
\end{theorem}

\begin{proof}
Let $A\colon V_\sigma\to V_\rho$ be linear, and average its translates:
\[
\widetilde A=\int_G\rho(g)A\sigma(g)^{-1}\,\dd\mu(g).
\]
These are integrals of matrix entries. Substitution $g'=hg$ shows that
$\rho(h)\widetilde A\sigma(h)^{-1}=\widetilde A$.
So $\widetilde A$ intertwines $\sigma$ and $\rho$. Schur's lemma makes it zero for distinct representatives. If $\rho=\sigma$, it is a scalar matrix, and taking its trace gives
$\widetilde A=(\operatorname{tr}A/d_\rho)I$.

To recover individual coefficients, let $A=E_{jl}$, the matrix with its only nonzero entry, equal to $1$, in row $j$ and column $l$. Its averaged entry in row $i$, column $k$ is
\[
\widetilde A_{ik}
=\int_G\rho_{ij}(g)\sigma(g^{-1})_{lk}\,\dd\mu(g)
=\int_G\rho_{ij}(g)\overline{\sigma_{kl}(g)}\,\dd\mu(g).
\]
When the two representations agree, the scalar-matrix formula gives $\delta_{jl}\delta_{ik}/d_\rho$. When they differ, it gives zero.

Sum the coefficient formula with $j=i$ and $l=k$. The result is
$\langle\chi_\rho,\chi_\sigma\rangle=\delta_{\rho\sigma}$.
By Proposition~\ref{prop:unitarytrick}, any character has a finite decomposition
$\chi=\sum_\rho m_\rho\chi_\rho$, where $m_\rho$ is the multiplicity of $\rho$. Orthogonality recovers it as $\langle\chi,\chi_\rho\rangle$. Equal characters have equal multiplicities, hence isomorphic representations.
\end{proof}

For $U(1)$, every irreducible is $z\mapsto z^k$, $k\in\Z$. Indeed, the commuting unitary matrices of a representation have a common eigenbasis, so an irreducible has dimension one; Volume~8, Theorem~13.11 classifies the continuous circle characters. Orthogonality becomes
\[
\frac1{2\pi}\int_0^{2\pi}e^{ik\theta}e^{-il\theta}\,\dd\theta=\delta_{kl}.
\]
This says that the Fourier modes are perpendicular. To expand every square-integrable function, we still need to know that no nonzero function is perpendicular to them all.

\section{The Peter--Weyl theorem}

Let $\widehat G$ denote our chosen representatives. Write $\mathcal R(G)$ for the complex linear span of matrix coefficients of all finite-dimensional representations. A sum in this span is finite. The theorem concerns its closure, where limits enter.

\begin{theorem}[Peter--Weyl, 1927]\label{thm:peterweyl}
\begin{enumerate}[label=(\alph*)]
\item The representative functions $\mathcal R(G)$ are uniformly dense in $C(G)$.
\item The functions $\sqrt{d_\rho}\rho_{ij}$, for $\rho\in\widehat G$ and $1\le i,j\le d_\rho$, form an orthonormal basis of $L^2(G,\mu)$. With
\[
\widehat f(\rho)=\int_G f(g)\rho(g)^*\,\dd\mu(g),
\]
every $f\in L^2(G,\mu)$ has the expansion
\[
f=\sum_{\rho\in\widehat G}d_\rho\operatorname{tr}\bigl(\widehat f(\rho)\rho\bigr),
\qquad
\int_G|f|^2\,\dd\mu
=\sum_\rho d_\rho\operatorname{tr}\bigl(\widehat f(\rho)\widehat f(\rho)^*\bigr).
\]
The series for $f$ converges in $L^2$.
\item The irreducible characters form an orthonormal basis of the closed subspace of class functions. Here a class function in $L^2$ means an equivalence class satisfying $f(hgh^{-1})=f(g)$ for every fixed $h$ and almost every $g$.
\end{enumerate}
\end{theorem}

\begin{proof}
\emph{Uniform approximation.} Products of coefficients are coefficients:
$\rho_{ij}\sigma_{kl}$ is the $((i,k),(j,l))$ entry of $\rho\otimes\sigma$.
The constant function $1$ comes from the trivial representation. Complex conjugates are available from the representation $g\mapsto\overline{\rho(g)}$. Thus $\mathcal R(G)$ is a conjugation-stable subalgebra of $C(G)$ containing the constants.

It also separates points. Two different matrices in $G\subseteq\mathrm{GL}_N(\C)$ differ in some entry. That entry is a coefficient of the defining representation on $\C^N$. The complex Stone--Weierstrass theorem, Volume~7, Theorem~4.19(b), now gives (a). The matrix-group assumption is exactly what made separation immediate.

\emph{From continuous functions to $L^2$.} Haar measure here is a finite regular Borel measure on the compact metric space $G$. Given a Borel set $E$ and $\varepsilon>0$, regularity supplies a closed $F$ and open $O$ with
$F\subseteq E\subseteq O$ and $\mu(O\setminus F)<\varepsilon$.
There is a continuous $0\le\phi\le1$ equal to $1$ on $F$ and $0$ off $O$.
When $F$ and $O^c$ are both nonempty, one explicit choice is
\[
\phi(x)=\frac{d(x,O^c)}{d(x,O^c)+d(x,F)}.
\]
Its denominator never vanishes because the two disjoint closed sets are compact. If $F$ is empty use $\phi=0$; if $O^c$ is empty use $\phi=1$. In either case the same error estimate holds:
$\|\phi-1_E\|_2^2\le\mu(O\setminus F)<\varepsilon$.
Completed-measure sets differ from Borel sets only by null sets, so they cause no extra difficulty. Simple functions are dense in $L^2$ by Volume~7's integration theory. Approximating their indicators proves that $C(G)$ is dense too.

Uniform approximation implies $L^2$ approximation since $\mu(G)=1$. Part (a) therefore makes $\mathcal R(G)$ dense in $L^2$. Every representation splits into irreducibles. In a basis respecting that splitting, its coefficients are irreducible coefficients or zero; a change of basis takes linear combinations. Consequently the orthonormal family in (b), orthonormal by Theorem~\ref{thm:schurorth}, has dense span and is a Hilbert-space basis.

The coefficient of $f$ on $\sqrt{d_\rho}\rho_{ij}$ is
$\sqrt{d_\rho}\int f\overline{\rho_{ij}}\,\dd\mu
=\sqrt{d_\rho}\widehat f(\rho)_{ji}$.
Insert these coefficients into the Hilbert-space expansion and Parseval identity of Volume~8. Grouping the $d_\rho^2$ terms for each $\rho$ gives the trace formulas in (b). The order of the indices explains why the Fourier transform uses $\rho(g)^*$.

\emph{Conjugation-invariant functions.} Let $E_\rho$ be the span of the coefficients of $\rho$, and let $P_\rho$ be its orthogonal projection. For fixed $h\in G$, conjugation acts by the unitary operator
$(c_hf)(g)=f(h^{-1}gh)$.
The identity $\rho(h^{-1}gh)=\rho(h)^{-1}\rho(g)\rho(h)$ shows that $E_\rho$ is invariant under $c_h$ and its inverse. Hence $P_\rho$ commutes with $c_h$.

If $f$ is a class function, $P_\rho f$ is fixed by every $c_h$. Write it uniquely as $g\mapsto\operatorname{tr}(A\rho(g))$; uniqueness follows from coefficient orthogonality. Conjugating $g$ replaces $A$ by $\rho(h)A\rho(h)^{-1}$. Equality in $L^2$ and coefficient independence force
$\rho(h)A\rho(h)^{-1}=A$ for all $h$.
Schur's lemma gives $A=aI$. Thus $P_\rho f=a\chi_\rho$, with $a=\langle f,\chi_\rho\rangle$. Part (b) now expands $f$ in characters. Conversely, an $L^2$ limit of their linear combinations is fixed by every unitary $c_h$. This proves both completeness and closedness in (c).
\end{proof}

On the circle, this is Fourier completeness. On $SO(3)$, coefficient functions restricted to functions invariant under a point stabiliser give the corresponding approximation on the sphere. The full Peter--Weyl theorem for arbitrary compact groups is a further result, quoted here: without a given faithful matrix representation, one must first prove that finite-dimensional representations separate points. Convolution operators and the compact-operator spectral theorem supply that argument. It is not part of the proof just given.

\begin{think}
The representation $t\mapsto\begin{pmatrix}1&t\\0&1\end{pmatrix}$ of $(\R,+)$ has the invariant line $\C e_1$. Show directly that no other line is invariant, so there is no invariant complementary line. Which step in the averaging proof fails when Haar measure has infinite total mass? Noncompactness does not by itself force every representation to fail: the question is whether an invariant complement exists for the particular representation.
\end{think}

\section{The irreducible representations of \texorpdfstring{$U(n)$}{U(n)}}

Every unitary matrix is unitarily diagonalizable (Volume~9, Theorem~7.9). A character is therefore determined by its values on the diagonal torus
$T=\{\operatorname{diag}(t_1,\ldots,t_n):|t_j|=1\}$.
Permutation matrices reorder the diagonal entries, so those values must be symmetric in the $t_j$. Weyl's integration formula, Volume~9, Theorem~8.8, says that for a continuous class function $F$,
\[
\int_{U(n)}F(U)\,\dd\mu(U)
=\frac1{n!}\int_T F(t)\prod_{j<k}|t_j-t_k|^2\,
\frac{\dd\theta_1\cdots\dd\theta_n}{(2\pi)^n},
\qquad t_j=e^{i\theta_j}.
\]
We will use this positive weight to turn symmetric polynomials into orthogonal characters.

For $\mu\in\Z^n$, set $t^\mu=\prod_jt_j^{\mu_j}$ and
\[
a_\mu(t)=\det(t_k^{\mu_j})_{j,k=1}^n
=\sum_{w\in\mathfrak S_n}\operatorname{sgn}(w)
 t_{w(1)}^{\mu_1}\cdots t_{w(n)}^{\mu_n}.
\]
This is an \emph{alternant}: swapping two variables changes its sign.
For $\delta=(n-1,n-2,\ldots,0)$, the Vandermonde identity from Volume~1, Proposition~7.4 gives
$a_\delta=\prod_{j<k}(t_j-t_k)$.
A \emph{dominant weight} is an integer vector $\lambda$ with
$\lambda_1\ge\cdots\ge\lambda_n$. Adding $\delta$ makes its entries strictly decreasing.

For $n=2$ and $\lambda=(m,0)$ with $m\ge0$, the quotient of alternants is already familiar:
\[
\frac{t_1^{m+1}-t_2^{m+1}}{t_1-t_2}
=t_1^m+t_1^{m-1}t_2+\cdots+t_2^m.
\]
The apparent division by zero at $t_1=t_2$ disappears. The next lemma proves that cancellation in every dimension.

\begin{definition}\label{def:schurpoly}
For a dominant weight $\lambda$, its Schur function is
$s_\lambda=a_{\lambda+\delta}/a_\delta$.
The quotient is first read where the coordinates are distinct, then extended by its Laurent-polynomial expression.
\end{definition}

\begin{lemma}\label{lem:alternants}
\begin{enumerate}[label=(\alph*)]
\item Each $s_\lambda$ is a symmetric Laurent polynomial with integer coefficients. In lexicographic order its largest monomial is $t^\lambda$, with coefficient $1$.
\item Every antisymmetric Laurent polynomial has a unique finite expression $\sum_\mu c_\mu a_\mu$ over strictly decreasing integer vectors $\mu$. Here $c_\mu$ is its coefficient of $t^\mu$.
\item If $\mu,\nu$ are strictly decreasing, then
\[
\frac1{n!}\int_T a_\mu(t)\overline{a_\nu(t)}\,
\frac{\dd\theta_1\cdots\dd\theta_n}{(2\pi)^n}
=\delta_{\mu\nu}.
\]
\end{enumerate}
\end{lemma}

\begin{proof}
\emph{Group monomials by their exponents.} In an antisymmetric Laurent polynomial, a monomial with two equal exponents has coefficient zero: exchange those two variables. For other exponents, exchanging variables changes the coefficient by the permutation's sign. Group each orbit by its unique decreasing exponent vector. It contributes precisely $c_\mu a_\mu$. Distinct decreasing vectors have disjoint monomial orbits, proving uniqueness and (b).

\emph{Cancel the Vandermonde.} Multiply the numerator by a sufficiently large power of $t_1\cdots t_n$ to make it a polynomial. Setting $t_j=t_k$ makes two columns equal, so $t_j-t_k$ divides it. These factors are pairwise nonassociate irreducibles in the unique factorisation domain $\Z[t_1,\ldots,t_n]$. Their product therefore divides it over $\Z$, not merely over $\Q$. Undo the initial monomial multiplication. The resulting quotient is an integer Laurent polynomial. Both numerator and denominator change sign under a variable transposition, so the quotient is symmetric.

Lexicographic order compares exponent vectors at the first coordinate where they differ. It is compatible with addition. The leading monomial of $a_\mu$ for decreasing $\mu$ is $t^\mu$, with coefficient $1$. Leading monomials multiply, so $a_{\lambda+\delta}=s_\lambda a_\delta$ forces the leading monomial of $s_\lambda$ to be $t^\lambda$, with coefficient $1$. This proves (a).

\emph{Integrate monomials.} A circle integral kills $z^r$ unless $r=0$. Iterating this gives $\int_T t^\alpha\overline{t^\beta}=\delta_{\alpha\beta}$. Each decreasing $\mu$ produces $n!$ distinct monomials with coefficients $\pm1$, and two different decreasing vectors produce disjoint sets. Dividing the resulting inner products by $n!$ gives (c).
\end{proof}

The integer coefficients matter. The norm calculation below will express $1$ as a sum of squares of integers, leaving exactly one surviving alternant.

\begin{theorem}[Weyl's character formula for $U(n)$, 1925]\label{thm:weylchar}
For each dominant weight $\lambda$, there is an irreducible representation $V_\lambda$ of $U(n)$ with character $s_\lambda$ on $T$. Every irreducible representation occurs for exactly one dominant weight. Its dimension is
\[
\dim V_\lambda=s_\lambda(1,\ldots,1)
=\prod_{1\le j<k\le n}
\frac{\lambda_j-\lambda_k+k-j}{k-j}.
\]
\end{theorem}

\begin{proof}
\emph{The weights of a representation.} The commuting unitary operators $\rho(t)$ have a common orthonormal eigenbasis. To see why finitely many suffice, diagonalize one operator; every operator commuting with it preserves its eigenspaces. Repeat on any eigenspace where some operator is not scalar. Each refinement lowers a positive dimension, so the process terminates.

On a common eigenline, $T$ acts by a continuous character. Its restriction to each circle factor is an integer power, by Volume~8, Theorem~13.11. Thus the character of $\rho$ is a symmetric Laurent polynomial
$\chi_\rho(t)=\sum_\mu m_\mu t^\mu$, where the nonnegative integer $m_\mu$ is the dimension of the weight space with eigenvalue $t^\mu$.

\emph{An irreducible character has only one alternant.} Multiply an irreducible character by $a_\delta$ and use Lemma~\ref{lem:alternants}(b):
$\chi_\rho a_\delta=\sum_\nu c_\nu a_\nu$, with integers $c_\nu$ and strictly decreasing $\nu$. Schur orthogonality and Weyl's integration formula give
\[
1=\int_{U(n)}|\chi_\rho|^2\,\dd\mu
=\frac1{n!}\int_T|\chi_\rho a_\delta|^2\,
\frac{\dd\theta}{(2\pi)^n}
=\sum_\nu c_\nu^2.
\]
Exactly one coefficient is $1$ or $-1$. Put $\lambda=\nu-\delta$ for its index; strict decrease of the integer entries of $\nu$ makes $\lambda$ dominant. We have $\chi_\rho=\pm s_\lambda$. The leading coefficient of $s_\lambda$ is $1$, whereas all monomial coefficients of $\chi_\rho$ are weight multiplicities. They cannot be negative. Hence the sign is positive. Distinct irreducibles have orthogonal characters, so no two give the same $\lambda$.

\emph{No dominant weight is missing.} A symmetric Laurent polynomial is a power of $t_1\cdots t_n$ times an ordinary symmetric polynomial. The latter is a polynomial in the elementary symmetric functions, by Volume~1, Theorem~6.12. On unitary matrices these elementary functions are coefficients of the characteristic polynomial, and $t_1\cdots t_n=\det U$ never vanishes. Therefore $s_\lambda$ defines a continuous class function of $U$, including repeated eigenvalues.

If it were not an irreducible character, the preceding step and alternant orthogonality would make it perpendicular to every irreducible character. Yet Weyl's integration formula gives $\|s_\lambda\|_2^2=1$. This contradicts the completeness of class characters in Theorem~\ref{thm:peterweyl}(c). Every dominant weight therefore occurs.

\emph{Evaluate at the identity.} Set $\nu=\lambda+\delta$ and approach $(1,\ldots,1)$ along
$t=(q^{n-1},q^{n-2},\ldots,1)$, with positive real $q\to1$. A Laurent polynomial can be evaluated along this path even though the path leaves the unit torus. A Vandermonde determinant gives
\[
a_\nu(q^{n-1},\ldots,1)
=\prod_{j<k}(q^{\nu_j}-q^{\nu_k}).
\]
Divide by the same product for $\delta$. Since
$(q^a-q^b)/(q-1)\to a-b$ for all integers $a,b$, the limit is
$\prod_{j<k}(\nu_j-\nu_k)/(\delta_j-\delta_k)$, the claimed formula. For $n=1$, the product is empty and has value $1$, as it should for a circle character.
\end{proof}

\begin{figure}[ht]
\centering
\fig{ch11-weights}
\caption{Weight diagrams of two representations of $U(3)$, drawn in the plane $\mu_1+\mu_2+\mu_3=\text{constant}$. Left: $V_{(2,1,0)}$, of dimension $8$: the six permutations of $(2,1,0)$, and the weight $(1,1,1)$ with multiplicity $2$ (the circled centre); restricted to $SU(3)$,
this is Gell-Mann's ``eightfold way'' of 1961, in which the eight lightest baryons fill the weights. Right: $V_{(3,0,0)}=\operatorname{Sym}^3\C^3$, of dimension $10$, whose tenth weight, at a corner of the triangle, was a particle not yet observed; Gell-Mann
predicted it in 1962, and the $\Omega^-$ was found at Brookhaven in 1964.}
\label{fig:weights}
\end{figure}

The two dimensions in Figure~\ref{fig:weights} follow directly from the product: for $(2,1,0)$ it gives $2\cdot2\cdot2=8$; for $(3,0,0)$ it gives $4\cdot(5/2)\cdot1=10$. A weight can occur more than once; the central weight of the first representation has multiplicity two. Counting points without their multiplicities would give the wrong dimension.

Several familiar representations now have names. The standard representation has character $p_1=t_1+\cdots+t_n=s_{(1,0,\ldots,0)}$. The exterior power $\Lambda^k\C^n$, $0\le k\le n$, has character $e_k$, the sum of squarefree degree-$k$ monomials. Its alternant product is
$e_k a_\delta=\sum_{|S|=k}a_{\delta+1_S}$.
Every term vanishes through repeated exponents unless $S=\{1,\ldots,k\}$. Thus $e_k=s_{(1^k,0^{n-k})}$, proving that the exterior power is irreducible.

The symmetric power $\operatorname{Sym}^k\C^n$, $k\ge0$, has as a basis the degree-$k$ monomials, so its character is the complete symmetric polynomial $h_k$. For distinct $t_i$, partial fractions of $\prod_i(1-t_i z)^{-1}$ give
\[
h_k=\sum_{i=1}^n\frac{t_i^{k+n-1}}{\prod_{j\ne i}(t_i-t_j)}
=\frac{a_{(k+n-1,n-2,\ldots,0)}}{a_\delta}.
\]
The last equality is expansion of the numerator determinant along its first row. Both sides are polynomials, so equality on distinct coordinates gives equality everywhere. This identifies $h_k=s_{(k,0,\ldots,0)}$ and proves irreducibility of the symmetric power too.

For $\lambda=(m,\ldots,m)$, the character is $(t_1\cdots t_n)^m$, so $V_\lambda=\det^m$. Tensoring with this one-dimensional representation adds $m$ to every entry of a weight. Any dominant weight can consequently be shifted into a partition, with all entries nonnegative. For a partition there is also the hook-content identity
$\dim V_\lambda=\prod_{(i,j)\in\lambda}(n+j-i)/h_{ij}$.
This alternative product is quoted here; the pairwise product already proved is enough for every dimension calculation below.

The actions of $\mathfrak S_k$ and $U(n)$ on $(\C^n)^{\otimes k}$ commute: the first permutes tensor factors and the second acts on each factor. Their full joint decomposition is a further theorem.

\begin{remark}[Schur--Weyl duality; quoted here]
The joint representation has the decomposition
\[
(\C^n)^{\otimes k}
\cong\bigoplus_{\lambda\vdash k,\ \ell(\lambda)\le n}S^\lambda\otimes V_\lambda.
\]
Schur's 1901 thesis is a source of this duality. We quote the joint isomorphism here and prove it in Chapter~14. The argument below proves its character consequence for $U(n)$ directly: the multiplicity of $V_\lambda$ is $f^\lambda=\dim S^\lambda$.
\end{remark}

\begin{figure}[ht]
\centering
\fig{v10-new-ch11-growth}
\caption{Adding boxes with at most two rows. There is one path to $(3)$ and two paths to $(2,1)$; the omitted shape $(1,1,1)$ has three rows. For $U(2)$ these multiplicities give the sixth absolute trace moment $1^2+2^2=5$.}
\label{fig:v10new-ch11-growth}
\end{figure}

\begin{figure}[ht]
\centering
\fig{v10-new-ch11-variance}
\caption{Exact power-trace variances for $U(4)$. The values $\min(k,4)$ rise to $4$ and then stay there. Four independent uniform points on the circle instead give variance $4$ at every positive integer frequency. The curves join integer values to guide the eye.}
\label{fig:v10new-ch11-variance}
\end{figure}

\section{Traces of random unitary matrices}

Let $U$ have Haar probability distribution on $U(n)$. To draw it at random means to integrate a function of $U$ against $\mu$; it does not mean drawing its entries independently. Its eigenvalues lie on the circle, with the joint density from Volume~9, Chapter~8. We can measure their dependence by adding their $k$th powers.

\begin{proposition}\label{prop:powertraces}
For positive integers $j,k$,
\[
\mathbb E[\operatorname{tr}U^j]=0,
\qquad
\mathbb E\bigl[\operatorname{tr}U^j\overline{\operatorname{tr}U^k}\bigr]
=\delta_{jk}\min(k,n).
\]
\end{proposition}

\begin{proof}
\emph{Rotate the whole spectrum.} If $|z|=1$, left invariance makes $zU$ Haar distributed. The mean of $\operatorname{tr}U^j$ is therefore $z^j$ times itself for every $z$, so it is zero. Likewise the displayed mixed expectation is $z^{j-k}$ times itself. It must vanish unless $j=k$.

\emph{Count the surviving alternants.} Write $p_k(t)=\sum_i t_i^k$. Symmetry of $p_k$ lets us expand
\[
p_k a_\delta=\sum_{i=1}^n a_{\delta+ke_i}.
\]
In the $i$th summand, the entry $n-i$ has been replaced by $n-i+k$. If $k<i$, the new entry duplicates the entry at position $i-k$, so the alternant is zero. If $k\ge i$, the new entry is at least $n$, above all the other entries. Reordering it produces, up to sign, an alternant with strictly decreasing exponents.

There are $\min(k,n)$ surviving indices. They give distinct decreasing exponent vectors: each removes a different entry $n-i$ from $\{0,\ldots,n-1\}$. The corresponding alternants are perpendicular, each of squared norm $n!$ on $T$. Weyl's integration formula thus yields
\[
\mathbb E|\operatorname{tr}U^k|^2
=\frac1{n!}\int_T|p_k a_\delta|^2\frac{\dd\theta}{(2\pi)^n}
=\min(k,n).
\]
\end{proof}

If $z_1,\ldots,z_n$ were independent uniform circle points, the same calculation would give $\mathbb E|\sum_i z_i^k|^2=n$ for every $k\ge1$: off-diagonal terms average to zero. Figure~\ref{fig:v10new-ch11-variance} shows the difference. At low frequencies the unitary eigenvalues cancel much more closely. In particular, $\operatorname{tr}U$ has variance $1$, whatever the positive size $n$.

Diaconis and Shahshahani proved that any fixed list of power traces tends jointly to independent complex Gaussians, with variances $1,2,\ldots$. We quote that joint limit. To prove the one-dimensional limit needed here, we will compute every moment of the first trace exactly. A small rule about adding boxes does the work.

\begin{lemma}[Pieri's rule for one box]\label{lem:pieri}
If $\lambda$ is dominant, then
$p_1s_\lambda=\sum_\nu s_\nu$, where $\nu$ runs over dominant weights obtained by adding $1$ to one entry of $\lambda$.
\end{lemma}

\begin{proof}
Multiplication by $p_1$ raises one exponent:
$p_1a_{\lambda+\delta}=\sum_i a_{\lambda+\delta+e_i}$.
For $i>1$, if $\lambda_{i-1}=\lambda_i$, the raised exponent equals the one immediately before it. That determinant vanishes. Otherwise $\lambda+e_i$ is still dominant, and no reordering or sign change is needed. The $i=1$ term always survives. Divide by $a_\delta$ on distinct coordinates, then use polynomial continuity to include the remaining coordinates.
\end{proof}

For a partition, a legal increase adds one box while leaving a Young diagram. Repeated multiplication by $p_1$ therefore builds diagrams one box at a time. Figure~\ref{fig:v10new-ch11-growth} shows the first three steps with the two-row restriction.

\begin{theorem}[Diaconis--Shahshahani 1994, Rains 1998]\label{thm:traces}
For each integer $a\ge0$,
\[
(\operatorname{tr}U)^a
=\sum_{\substack{\lambda\vdash a\\\ell(\lambda)\le n}}
 f^\lambda\chi_{V_\lambda}(U),
\]
and
\[
\mathbb E|\operatorname{tr}U|^{2a}
=\#\{\sigma\in\mathfrak S_a:
 \text{the longest increasing subsequence of $\sigma$ has length at most $n$}\}.
\]
In particular, this moment is $a!$ whenever $a\le n$.
\end{theorem}

\begin{proof}
Start with the empty partition and the constant character $1$. Apply Pieri's rule $a$ times. The coefficient of $s_\lambda$ is the number of paths to $\lambda$ that add one box at each step and never exceed $n$ rows.

Label each new box by its time of arrival. Boxes to its left or above must already have arrived. The result is a standard tableau. Conversely, remove the largest label from a standard tableau, then the next largest, and continue: every intermediate shape is a Young diagram. These operations are inverse, so the coefficient is $f^\lambda$. The character equality holds on $T$, hence on $U(n)$ by diagonalization.

Take the squared $L^2$ norm. Character orthogonality gives
\[
\mathbb E|\operatorname{tr}U|^{2a}
=\sum_{\lambda\vdash a,\ \ell(\lambda)\le n}(f^\lambda)^2.
\]
Transposing tableaux exchanges rows and columns and preserves their number (Corollary~\ref{cor:conjugate}). We may replace the row bound by $\lambda_1\le n$. Robinson--Schensted pairs standard tableaux of equal shape with permutations, and Schensted's theorem identifies the first row length with the longest increasing subsequence (Theorems~\ref{thm:rs} and \ref{thm:schensted}). The sum is therefore exactly the stated count.

When $a\le n$, every permutation qualifies, giving $a!$. When $a=0$, there is one empty tableau and one empty permutation; the integral of $1$ is also $1$. This includes the boundary case.
\end{proof}

For $U(2)$ and $a=3$, the only excluded permutation is $123$. The moment is $5$, just as Figure~\ref{fig:v10new-ch11-growth} predicts. This is a finite equality, before any limit is taken.

Gessel's 1990 generating function expresses the same permutation counts as a Toeplitz determinant. Baik, Deift and Johansson used that identity in their 1999 analysis of the longest increasing subsequence. These later results explain why a question about random permutations also leads to integrals over unitary matrices; we do not need either result for the moment identity.

\begin{corollary}\label{cor:gaussiantrace}
As $n\to\infty$, $\sqrt2\operatorname{Re}\operatorname{tr}U$ converges in distribution to $N(0,1)$. For every fixed real $\phi$, the same is true of $\sqrt2\operatorname{Re}(e^{-i\phi}\operatorname{tr}U)$.
\end{corollary}

\begin{proof}
Set $X_n=\operatorname{tr}U$. Rotational invariance gives
$\mathbb E[X_n^j\overline{X_n}^l]=0$ whenever $j\ne l$.
Expand $\operatorname{Re}X_n=(X_n+\overline{X_n})/2$. Every odd moment vanishes. For an even moment $2r$, once $n\ge r$, only the middle binomial term remains, and Theorem~\ref{thm:traces} gives
\[
\mathbb E[(\operatorname{Re}X_n)^{2r}]
=2^{-2r}\binom{2r}{r}r!
=\frac{(2r)!}{4^r r!}.
\]
Multiplying by $2^r$ yields $(2r)!/(2^r r!)$, the $2r$th standard Gaussian moment. Thus every fixed moment of $\sqrt2\operatorname{Re}X_n$ eventually equals its Gaussian counterpart. Volume~6, Lemma~9.6 proves that convergence of all moments to these Gaussian moments implies convergence in distribution; its proof uses Taylor's formula and Gaussian smoothing, with no premise from Volume~11. Apply it here. Finally $e^{-i\phi}X_n$ has the same distribution as $X_n$, so the rotated assertion follows as well.
\end{proof}

The trace has fluctuations of size $1$, rather than $\sqrt n$. Its limiting law is a standard complex Gaussian, meaning that its real and imaginary parts tend jointly to independent $N(0,1/2)$ variables; all real projections follow from the corollary, and the joint formulation is the usual projection criterion, quoted here. For a fixed finite $n$ it is not exactly Gaussian: $|\operatorname{tr}U|\le n$, while a Gaussian has unbounded support. Exact agreement of each fixed low-order moment for sufficiently large $n$ is compatible with that difference.

\begin{beauty}
Keating and Snaith used Weyl integration in 2000 to compute moments of $|\det(I-U)|$. After an arithmetic factor is included, their unitary-matrix calculation predicts constants in conjectures for high moments of $\zeta(\tfrac12+it)$. The predictions agree with the known second- and fourth-moment leading terms, due to Hardy--Littlewood and Ingham. At the sixth and eighth moments they yield the factors $42$ and $24024$ also found in arithmetic heuristics. These higher zeta-moment predictions remain conjectural; agreement of constants is evidence, not a proof about zeta zeros. Exercise~\ref{ex:charpoly} proves the simplest unitary calculation, $\mathbb E|\det(I-U)|^2=n+1$.
\end{beauty}

% V10-EXPANSION WORKED EXAMPLE 1
\subsection*{Worked example: a negative weight still gives a finite representation}

On $U(2)$, take the dominant weight $(2,-1)$. Subtracting the final entry from both entries leaves $(3,0)$, so its representation is
\[
 V_{(2,-1)}\cong\det^{-1}\otimes\operatorname{Sym}^3\C^2.
\]
This uses only the determinant twist and symmetric-power character already proved. The character is the Laurent polynomial
\[
 s_{(2,-1)}(t_1,t_2)
 =\frac{t_1^2}{t_2}+t_1+t_2+\frac{t_2^2}{t_1}.
\]
At the identity the value is $4$, matching the dimension product $2-(-1)+1=4$. The negative exponent does not create an infinite-dimensional space; it describes the action of a one-dimensional determinant inverse on an ordinary four-dimensional space.

A central matrix $zI$ acts by $z$, since every displayed monomial has total degree $1$. On $SU(2)$, where $t_1=e^{i\alpha}$ and $t_2=e^{-i\alpha}$, the same character becomes $2\cos3\alpha+2\cos\alpha$. At $\alpha=0$ it is $4$, at $\pi$ it is $-4$, and at $\pi/2$ it is zero. These values check the signs of the negative powers.

Tensor with the standard representation. Pieri's one-box rule gives
\[
 V_{(2,-1)}\otimes\C^2\cong V_{(3,-1)}\oplus V_{(2,0)}.
\]
Both increases preserve dominance. Their dimensions are $5$ and $3$, adding to $8=4\cdot2$, and both have central exponent $2$, as tensor multiplication requires. This is a character calculation followed by complete reducibility; it does not use the later joint Schur--Weyl decomposition.

% V10-EXPANSION WORKED EXAMPLE 2
\subsection*{Worked example: where the eighth trace moment stops being Gaussian}

Let $X_n=\operatorname{tr}U$ for Haar $U\in U(n)$. For $n=2$, the fourth power of its character has three allowed shapes:
\[
 p_1^4=s_{(4)}+3s_{(3,1)}+2s_{(2,2)}.
\]
These coefficients count standard tableaux. The corresponding $U(2)$ dimensions are $5,3,1$, so evaluation at the identity checks $16=5+3\cdot3+2\cdot1$. Taking the squared character norm instead gives
\[
 \mathbb E|X_2|^8=1^2+3^2+2^2=14.
\]
Dimensions and multiplicities are doing different jobs in these two equalities.

At size $3$, the shape $(2,1,1)$ is also allowed and contributes $3^2=9$, giving moment $23$. At size $4$, the final shape $(1,1,1,1)$ contributes $1$, giving $24$. The first four radial moments are therefore
\[
\begin{array}{c|rrrr}
 &\mathbb E|X_n|^2&\mathbb E|X_n|^4&\mathbb E|X_n|^6&\mathbb E|X_n|^8\\\hline
 n=1&1&1&1&1\\
 n=2&1&2&5&14\\
 n=3&1&2&6&23\\
 n\ge4&1&2&6&24
\end{array}
\]
A standard complex Gaussian has this final row, since its radial moment of order $2a$ is $a!$. The exact-match range for $U(n)$ is $a\le n$, not all moments at once.

Rotation also turns a radial moment into a real moment. In $(X_n+\overline{X_n})^8$, only the balanced term survives integration. Consequently
\[
 \mathbb E(\sqrt2\operatorname{Re}X_2)^8
   =\frac1{16}\binom84\,14=\frac{245}{4},
\]
whereas the eighth standard Gaussian moment is $105$. This finite discrepancy does not oppose the limiting theorem. For a fixed eighth moment, the equality becomes exact once $n\ge4$; the limit lets $n$ grow while the moment order stays fixed.

\section{What we proved, and what comes next}

An invariant average made every finite-dimensional representation unitary. Stone--Weierstrass then made its coefficient functions complete. On $U(n)$ the Vandermonde weight in Weyl integration turned character orthogonality into an integer norm calculation. That forced the irreducible characters to be Schur functions. Adding boxes to those functions converted a trace integral into a count of permutations.

The proofs here establish Peter--Weyl for compact matrix groups, Weyl's character and dimension formulas for $U(n)$, the power-trace covariances and the permutation moment identity. The general compact-group extension, hook-content identity, joint power-trace limit and higher zeta-moment predictions were identified separately as quoted results or conjectures. Full Schur--Weyl duality is quoted here and proved in Chapter~14; no proof in this chapter depends on that later argument.

Three concluding chapters put the algebra to work on computations that deserve more room. Chapter~12 returns to projective resolutions and calculates their homology maps explicitly. Chapter~13 works through group extensions and cohomology. Chapter~14 returns to the two commuting actions on tensor powers and proves the joint Schur--Weyl decomposition quoted here. A reader following the representation route can continue directly to that last chapter; the next printed chapter takes up the resolution calculations. Exact finite trace integrals will remain our starting point when the later probability volumes ask about limiting laws.

\section{Exercises}

\begin{exercise}[$\star$]\label{ex:dims}
For $n\ge1$ and integer $k\ge0$, compute $\dim\operatorname{Sym}^k\C^n$ with Theorem~\ref{thm:weylchar}. Compute $\dim\Lambda^k\C^n$ for $0\le k\le n$. For $n\ge2$, compute the dimension of $V_{(1,0,\ldots,0,-1)}$ and identify it in $\operatorname{End}(\C^n)\cong\C^n\otimes\overline{\C^n}$, with the conjugation action.
\end{exercise}

\begin{exercise}[$\star$]\label{ex:tensorsquare}
Take $n\ge2$. Prove
$s_{(1,0,\ldots,0)}^2=s_{(2,0,\ldots,0)}+s_{(1,1,0,\ldots,0)}$.
Use it to identify the two irreducible summands in
$\C^n\otimes\C^n=\operatorname{Sym}^2\C^n\oplus\Lambda^2\C^n$.
\end{exercise}

\begin{exercise}[$\star\star$]\label{ex:charpoly}
Express $\det(I-U)$ as $\sum_{k=0}^n(-1)^k\chi_{\Lambda^k}(U)$. Use orthogonality to prove $\mathbb E|\det(I-U)|^2=n+1$ for Haar $U\in U(n)$.
\end{exercise}

\begin{exercise}[$\star\star$]\label{ex:su2}
Let $m$ be a nonnegative integer. Restrict $V_{(m,0)}$ to $SU(2)$. Show that its character at eigenvalues $e^{\pm i\alpha}$ is
$\chi_m(\alpha)=\sin((m+1)\alpha)/\sin\alpha$, with continuous values at $0$ and $\pi$. Use the density $(2/\pi)\sin^2\alpha$ on $[0,\pi]$ from Volume~9, Exercise~8.3 to prove these representations are irreducible, pairwise non-isomorphic, and exhaust the irreducible representations of $SU(2)$.
\end{exercise}

\begin{exercise}[$\star\star$]\label{ex:catalanu2}
For Haar $U\in U(2)$ and integer $a\ge0$, evaluate $\mathbb E|\operatorname{tr}U|^{2a}$ directly from Weyl integration. Combine it with Theorem~\ref{thm:traces} to count permutations of $\{1,\ldots,a\}$ with no increasing subsequence of length $3$.
\end{exercise}

\begin{exercise}[$\star\star\star$]\label{ex:domino}
Prove $p_2s_\lambda=\sum_\nu\pm s_\nu$. The sum is over dominant weights obtained by a domino addition: raise two adjacent equal row lengths by $1$ (a vertical domino, sign $-$), or raise one row length by $2$ (a horizontal domino, sign $+$), keeping the resulting vector dominant. For a general integer dominant weight, first shift every entry by the same integer to draw its partition diagram. This is the two-box case of the Murnaghan--Nakayama rule.
\end{exercise}

\begin{exercise}[$\star\star$]\label{ex:v10new-ch11-statistic}
Let $f(z)=\sum_{k=-M}^M c_kz^k$ be a Laurent polynomial and $U\in U(n)$ Haar. Define $L_f(U)=\sum_{j=1}^n f(t_j)$ from its eigenvalues. Find $\mathbb E L_f$ and $\mathbb E|L_f-\mathbb E L_f|^2$. Compare the answer with $n$ independent uniform circle points, and compute it for $f(e^{i\theta})=\cos\theta+2\cos2\theta$.
\end{exercise}

\begin{exercise}[$\star\star$]\label{ex:v10new-ch11-detmoment}
For integers $r$ and nonnegative integers $a,b$, prove
$\mathbb E[(\det U)^r(\operatorname{tr}U)^a\overline{\operatorname{tr}U}^b]=0$
unless $nr+a-b=0$. Decide whether this condition also guarantees a nonzero expectation, using the tableau expansion of trace powers. Compute $\mathbb E[\overline{\det U}(\operatorname{tr}U)^n]$ and explain the answer using the one-column tableau.
\end{exercise}

\begin{exercise}[$\star\star$]\label{ex:v10new-ch11-projective}
The projective unitary group is $PU(n)=U(n)/\{zI:|z|=1\}$. Prove that $V_\lambda$ factors through this quotient exactly when $\sum_j\lambda_j=0$. Given a partition $\mu$, determine when a determinant twist $\det^r\otimes V_\mu$ factors through $PU(n)$, and find $r$. Check the standard representation and the trace-zero conjugation representation for $n\ge2$.
\end{exercise}

\begin{exercise}[$\star\star$]\label{ex:v10new-ch11-firstfailure}
Write $X_n=\operatorname{tr}U$ for Haar $U\in U(n)$. Prove
$\mathbb E|X_n|^{2(n+1)}=(n+1)!-1$.
For $t>0$ and $1\le r\le n$, prove
$P(|X_n|\ge t)\le r!/t^{2r}$.
Explain why the exact lower moments and these tail bounds do not make $X_n$ exactly Gaussian at finite $n$.
\end{exercise}

\begin{exercise}[$\star$]\label{ex:v10exp-ch11-weight27}
For $U(3)$, compute the dimension of $V_{(3,1,-1)}$, rewrite it as a determinant twist of a representation indexed by a partition, and determine the action of every central scalar $zI$. Explain why the sum of the weight entries is not its dimension.
\end{exercise}

\begin{exercise}[$\star\star$]\label{ex:v10exp-ch11-matrixcov}
For fixed complex $n\times n$ matrices $A,B$ and Haar $U\in U(n)$, prove
\[
 \mathbb E\operatorname{tr}(AU)=0,\qquad
 \mathbb E[\operatorname{tr}(AU)\overline{\operatorname{tr}(BU)}]
     =\frac1n\operatorname{tr}(AB^*).
\]
Deduce that $n|\operatorname{tr}(AU)|^2$ is an unbiased estimate of $\operatorname{tr}(AA^*)$. Check a matrix unit and $A=I$ separately, and identify which orthogonality formula makes independence of matrix entries unnecessary.
\end{exercise}

\begin{exercise}[$\star\star$]\label{ex:v10exp-ch11-dependent}
For Haar $U\in U(2)$, put $P=\operatorname{tr}U$ and $Q=\operatorname{tr}U^2$. Show that $\mathbb E[P\overline Q]=0$, but $P,Q$ are not independent, by computing $\mathbb E|PQ|^2$. Prove the identity $p_1p_2=s_{(3,0)}$ in two variables. Compare with the value of $\mathbb E|PQ|^2$ on $U(3)$, without claiming that equality with a product of variances alone would prove independence.
\end{exercise}

\begin{exercise}[$\star\star\star$]\label{ex:v10exp-ch11-fourthpowertrace}
Starting from alternants, derive
\[
 p_2^2=s_{(4)}-s_{(3,1)}+2s_{(2,2)}-s_{(2,1,1)}+s_{(1^4)},
\]
with any shape having more than $n$ rows omitted for $U(n)$. Compute $\mathbb E|\operatorname{tr}U^2|^4$ for every positive $n$, and compare it with the fourth absolute moment of a centred complex Gaussian of variance $2$. Explain what the finite-size discrepancies do and do not say about the quoted joint Gaussian limit.
\end{exercise}

\section{Hints}

\begin{hint}{ex:dims}
Most factors of the dimension product equal $1$. Separate the pairs touching the first or last row. For the final identification, split scalar matrices from trace-zero matrices and inspect the largest torus weight.
\end{hint}

\begin{hint}{ex:tensorsquare}
Add one box to $(1)$ using Pieri's rule. The flip of two tensor factors has eigenspaces $\operatorname{Sym}^2$ and $\Lambda^2$.
\end{hint}

\begin{hint}{ex:charpoly}
Expand $\prod_j(1-t_j)$. The elementary symmetric functions are distinct irreducible characters, including the constant character at $k=0$.
\end{hint}

\begin{hint}{ex:su2}
Substitute $t_1=e^{i\alpha}$ and $t_2=e^{-i\alpha}$ into the two-variable alternant. Multiplication by $\sin\alpha$ turns the weighted character inner product into the ordinary sine-series inner product.
\end{hint}

\begin{hint}{ex:catalanu2}
Use $\varphi=\theta_1-\theta_2$. Then $|t_1+t_2|^2=4\cos^2(\varphi/2)$ and $|t_1-t_2|^2=4\sin^2(\varphi/2)$. Subtract two consecutive Wallis integrals.
\end{hint}

\begin{hint}{ex:domino}
Multiply the alternant by $p_2$, raising one exponent by $2$. A collision makes the term zero. Crossing one neighbouring exponent changes the sign.
\end{hint}

\begin{hint}{ex:v10new-ch11-statistic}
Put $p_k=\operatorname{tr}U^k$ for all nonzero integers $k$; then $p_{-k}=\overline{p_k}$. Rotation makes distinct signed frequencies perpendicular.
\end{hint}

\begin{hint}{ex:v10new-ch11-detmoment}
Replace $U$ by $zU$. For the requested value, take the character inner product of $p_1^n$ with $\det$ and apply the tableau expansion.
\end{hint}

\begin{hint}{ex:v10new-ch11-projective}
The central scalar $zI$ acts by a scalar on an irreducible representation. Evaluate its character using the homogeneity of the alternant quotient. A representation descends precisely when that scalar is always $1$.
\end{hint}

\begin{hint}{ex:v10new-ch11-firstfailure}
Among permutations of $n+1$ letters, how many have an increasing subsequence of length $n+1$? Apply Markov's inequality to the nonnegative variable $|X_n|^{2r}$ for the tail estimate.
\end{hint}

\begin{hint}{ex:v10exp-ch11-weight27}
Add $1$ to each entry to get the partition $(4,2,0)$. The three dimension factors are equal. Total degree governs a central scalar.
\end{hint}
\begin{hint}{ex:v10exp-ch11-matrixcov}
Expand $\operatorname{tr}(AU)=\sum_{i,j}A_{ij}U_{ji}$ and apply Schur orthogonality to the defining representation, whose dimension is $n$.
\end{hint}
\begin{hint}{ex:v10exp-ch11-dependent}
In two variables, $p_1p_2=t_1^3+t_1^2t_2+t_1t_2^2+t_2^3$. In three variables the correction is the exterior-cube character. Compare a product moment with the product of the separate variances.
\end{hint}
\begin{hint}{ex:v10exp-ch11-fourthpowertrace}
First obtain $p_2=s_{(2)}-s_{(1,1)}$. Raising an alternant exponent by two gives positive horizontal additions and negative vertical additions. Keep the two contributions to $(2,2)$ before combining coefficients.
\end{hint}

\section{Solutions}

\begin{solution}{ex:dims}
\idea Cancel factors equal to $1$, then telescope the remaining products. The largest weight identifies the conjugation representation.
\writeup For $\lambda=(k,0,\ldots,0)$ only pairs $(1,l)$ contribute. Their product is
\[
\prod_{l=2}^n\frac{k+l-1}{l-1}=\binom{n+k-1}{k}.
\]
This includes $k=0$, giving the trivial representation. For $\lambda=(1^k,0^{n-k})$, the nontrivial factors have $j\le k<l$. For fixed $j$ their product is $(n-j+1)/(k-j+1)$. Multiplying over $j$ gives $\binom nk$. Empty products handle $k=0$ and $k=n$.

For $\lambda=(1,0,\ldots,0,-1)$, $n\ge2$, the pairs $(1,l)$ with $l<n$ contribute $n-1$ together, as do the pairs $(j,n)$ with $j>1$. The pair $(1,n)$ contributes $(n+1)/(n-1)$. The product is $n^2-1$. For $n=2$ the two earlier products are empty and equal $1$, giving $3$ as required.

Conjugation on $\operatorname{End}(\C^n)$ has character
$|\sum_jt_j|^2=n+\sum_{j\ne k}t_jt_k^{-1}$.
The scalar line splits off, and the trace-zero complement has character
$n-1+\sum_{j\ne k}t_jt_k^{-1}$, whose leading monomial is $t_1t_n^{-1}$. Decompose this complement into irreducibles. Its leading monomial must be the largest leading monomial of a constituent, since weight multiplicities are nonnegative. Hence $V_{(1,0,\ldots,0,-1)}$ occurs. Its dimension already equals the complement's dimension, so it is the entire complement. The action is on $\mathfrak{sl}_n(\C)$, the trace-zero matrices.
\end{solution}

\begin{solution}{ex:tensorsquare}
\idea Pieri's two possible additions match the symmetric and alternating pieces.
\writeup Adding a box to $(1)$ gives $(2)$ or $(1,1)$, since $n\ge2$. Thus $p_1^2=s_{(2)}+s_{(1,1)}$. The tensor square has this character, so it is the direct sum of these two irreducibles.

The flip $v\otimes w\mapsto w\otimes v$ commutes with $U(n)$. Its $+1$ and $-1$ eigenspaces give the symmetric and exterior squares; the projections are half the identity plus or minus half the flip. Both spaces are nonzero. Since there are only two irreducible constituents in total, each is irreducible. The weight $t_1^2$, represented by $e_1\otimes e_1$, occurs in the symmetric square and not in the exterior square. It distinguishes $V_{(2)}$ from $V_{(1,1)}$ and makes the identification precise.
\end{solution}

\begin{solution}{ex:charpoly}
\idea An alternating sum of distinct orthonormal characters has squared norm equal to the number of terms.
\writeup Diagonalization gives
$\det(I-U)=\prod_j(1-t_j)=\sum_{k=0}^n(-1)^ke_k(t)$.
The function $e_k$ is the character of $\Lambda^k\C^n$, identified above as the irreducible of weight $(1^k,0^{n-k})$. These $n+1$ weights are distinct, including the two one-dimensional representations at $k=0,n$. Character orthogonality now kills every cross term and leaves
$\mathbb E|\det(I-U)|^2=\sum_{k=0}^n1=n+1$.
As a check, at $n=1$, $\mathbb E|1-z|^2=2$. The $k=1$ case of the quoted Keating--Snaith product also telescopes:
$\prod_{j=0}^{n-1}\Gamma(j+1)\Gamma(j+3)/\Gamma(j+2)^2
=\prod_{j=0}^{n-1}(j+2)/(j+1)=n+1$.
\end{solution}

\begin{solution}{ex:su2}
\idea The Vandermonde quotient becomes a sine quotient, and its apparent singularities cancel. Fourier completeness rules out missing irreducibles.
\writeup The two-variable formula gives
\[
s_{(m,0)}(e^{i\alpha},e^{-i\alpha})
=\frac{e^{i(m+1)\alpha}-e^{-i(m+1)\alpha}}
 {e^{i\alpha}-e^{-i\alpha}}
=\frac{\sin((m+1)\alpha)}{\sin\alpha}.
\]
At $\alpha=0$ the limit is $m+1$; at $\pi$ it is $(-1)^m(m+1)$. These also follow from the scalar matrices $I$ and $-I$ acting on the symmetric power.

Using the stated density, the inner product of $\chi_m$ and $\chi_l$ is
\[
\frac2\pi\int_0^\pi\sin((m+1)\alpha)\sin((l+1)\alpha)\,\dd\alpha
=\delta_{ml}.
\]
Complete reducibility says that a character's squared norm is the sum of the squared multiplicities of its irreducible constituents. Norm $1$ therefore makes each restriction irreducible. Their orthogonality makes them pairwise different.

If an irreducible character $g$ were missing, it would be perpendicular to every $\chi_m$. The continuous function $g(\alpha)\sin\alpha$ would be perpendicular to every sine on $[0,\pi]$. Its odd extension to $[-\pi,\pi]$ is continuous and periodic, because its endpoint values are zero. Every Fourier coefficient vanishes: sine coefficients by the orthogonality just stated and cosine coefficients by oddness. The circle case of Peter--Weyl makes that function zero in $L^2$, hence everywhere by continuity. Thus $g=0$ on $(0,\pi)$ and, by continuity, at the endpoints too. An irreducible character has norm $1$, a contradiction. No irreducible is missing.
\end{solution}

\begin{solution}{ex:catalanu2}
\idea The Vandermonde factor inserts a sine square, making the difference of two Wallis integrals.
\writeup The integrand depends only on $\varphi=\theta_1-\theta_2$. Periodicity reduces the normalized two-torus integral to a normalized circle integral. With $x=\varphi/2$, Weyl's factor $1/2!$ gives
\[
\mathbb E|\operatorname{tr}U|^{2a}
=\frac{2\cdot4^a}\pi\int_0^\pi\cos^{2a}x(1-\cos^2x)\,\dd x.
\]
Let $I_a=\int_0^\pi\cos^{2a}x\,\dd x=\pi4^{-a}\binom{2a}{a}$, the Wallis integral already established in Volume~1's integration chapter. Substitution yields
\[
\frac{2\cdot4^a}\pi(I_a-I_{a+1})
=2\binom{2a}{a}-\frac12\binom{2a+2}{a+1}
=\frac1{a+1}\binom{2a}{a}.
\]
The last equality uses $\binom{2a+2}{a+1}=2(2a+1)\binom{2a}{a}/(a+1)$. It also works for $a=0$, giving $1$. Theorem~\ref{thm:traces} identifies the integral with the desired permutation count. This is the Catalan number $C_a$.
\end{solution}

\begin{solution}{ex:domino}
\idea Raising an alternant exponent by two can cross at most one distinct integer exponent without either stopping or colliding. That crossing accounts for the negative sign.
\writeup Put $l=\lambda+\delta$, so $l_1>\cdots>l_n$, and expand
$p_2a_l=\sum_i a_{l+2e_i}$.
For $i=1$, or when $l_{i-1}\ge l_i+3$, the raised vector is still decreasing. It contributes $s_{\lambda+2e_i}$ after division by $a_\delta$: a horizontal domino with positive sign.

If $l_{i-1}=l_i+2$, two entries coincide and the determinant is zero. If $l_{i-1}=l_i+1$, the raised entry crosses the previous one. Swapping them introduces a minus sign. When $i\ge3$ and $l_{i-2}=l_i+2$, there is a new collision, so again the determinant is zero. Otherwise the reordered vector is strictly decreasing; the weight has increased by $1$ in each of rows $i-1,i$. The condition $l_{i-1}=l_i+1$ means those rows had equal lengths. The absence of the second collision is exactly the condition that their vertical addition leaves a diagram.

Every legal horizontal or vertical domino determines its unique raised index $i$, so no term is lost or counted twice. This proves the signed sum on distinct coordinates; polynomial continuation covers repeated coordinates. Shifting all entries of $\lambda$ multiplies every term by the same determinant power and leaves the argument unchanged.
\end{solution}

\begin{solution}{ex:v10new-ch11-statistic}
\idea Traces are Fourier modes of the eigenvalue count. Distinct signed frequencies are perpendicular, though their variances depend on the matrix size.
\writeup Set $p_k=\sum_jt_j^k$ for $k\ne0$. Then
$L_f=nc_0+\sum_{k\ne0}c_kp_k$ and $\mathbb E p_k=0$, also for negative $k$ by conjugation. For signed nonzero integers $k,l$, rotational invariance gives
$\mathbb E[p_k\overline{p_l}]=0$ if $k\ne l$.
For $k=l$, Proposition~\ref{prop:powertraces} and conjugation give $\mathbb E|p_k|^2=\min(|k|,n)$. Therefore
\[
\mathbb E L_f=nc_0,\qquad
\mathbb E|L_f-nc_0|^2=\sum_{k\ne0}\min(|k|,n)|c_k|^2.
\]
For independent uniform points, orthogonality of circle monomials instead gives $n\sum_{k\ne0}|c_k|^2$.

For $\cos\theta+2\cos2\theta$, the coefficients at $\pm1$ are $1/2$ and at $\pm2$ are $1$. The mean is zero and the unitary variance is $1/2+2\min(2,n)$: it is $5/2$ at $n=1$ and $9/2$ for $n\ge2$. The independent variance is $5n/2$. A constant polynomial has both variances zero, as the formula predicts.
\end{solution}

\begin{solution}{ex:v10new-ch11-detmoment}
\idea A central rotation records the total weight. Orthogonality then extracts the determinant's coefficient.
\writeup Replacing $U$ by $zU$ multiplies the integrand by $z^{nr+a-b}$. If this integer is nonzero, choose $z$ on the circle with that power different from $1$. Equality of the two expectations forces the expectation to vanish. For the converse, first take $r\ge0$ and $b=a+nr$. Twisting by $\det^r$ adds $r$ to every entry of a weight. Expand both trace powers in characters using Theorem~\ref{thm:traces}. Orthogonality gives the finite sum
\[
\sum_{\substack{\lambda\vdash a\\\ell(\lambda)\le n}}
f^\lambda f^{\lambda+(r^n)}.
\]
Here $\lambda$ is padded to $n$ rows, and the shifted diagram has $b$ boxes. Every summand is a positive integer: every finite partition diagram has a standard tableau, for example the row-by-row filling. There is at least one $\lambda$, including the empty diagram when $a=0$. Thus the expectation is positive. For $r<0$, take the complex conjugate and interchange $a,b$; this reduces to the case just proved. In this family the balance condition is both necessary and sufficient.

For the requested value, $\det$ is $\chi_{(1^n)}$. Theorem~\ref{thm:traces} expands $p_1^n$ in characters, and its coefficient at this weight is $f^{(1^n)}=1$: a column has just one standard filling. Character orthogonality gives
$\mathbb E[\overline{\det U}(\operatorname{tr}U)^n]=1$.
\end{solution}

\begin{solution}{ex:v10new-ch11-projective}
\idea A quotient representation exists exactly when the subgroup being collapsed acts trivially.
\writeup Let $s=\sum_j\lambda_j$. Every monomial of $s_\lambda$ has total degree $s$, because the numerator and denominator alternants have degrees differing by $s$. Hence
$\chi_\lambda(zI)=z^s\dim V_\lambda$.
The central operator $\rho(zI)$ is scalar by Schur's lemma, so this trace equality makes it $z^s I$. It equals $I$ for every unit complex $z$ exactly when $s=0$.

If this holds, define the quotient representation by $[U]\mapsto\rho(U)$. It is well defined and continuous by the definition of the quotient topology. Conversely any quotient representation must kill every $zI$, proving the criterion. The determinant twist changes the sum to $|\mu|+nr$. An integer $r$ exists exactly when $n$ divides $|\mu|$, and it is uniquely $-|\mu|/n$. For $n\ge2$, the standard weight has sum $1$ and does not descend. The trace-zero conjugation weight $(1,0,\ldots,0,-1)$ has sum zero and does descend. For $n=1$, the quotient is trivial, consistent with the zero-weight condition.
\end{solution}

\begin{solution}{ex:v10new-ch11-firstfailure}
\idea The first missing tableau count measures the first missing Gaussian radial moment. A low moment still gives a valid tail bound.
\writeup A permutation of $n+1$ letters has an increasing subsequence of length $n+1$ exactly when the whole permutation is increasing. There is one such permutation. Theorem~\ref{thm:traces} therefore gives $(n+1)!-1$. At $n=1$, this reads $\mathbb E|U|^4=1=2!-1$, a useful boundary check.

For $1\le r\le n$, the exact moment is $\mathbb E|X_n|^{2r}=r!$. Markov's inequality applied to $|X_n|^{2r}$ gives $P(|X_n|\ge t)\le r!/t^{2r}$. The bound can be replaced by its minimum with $1$; it need not be sharp.

A standard complex Gaussian $Z$, with planar density $\pi^{-1}e^{-|z|^2}$, has $\mathbb E|Z|^{2a}=a!$ by polar integration and the gamma integral. The first displayed moment differs from that Gaussian value by $1$. Moreover $|X_n|\le n$ while $Z$ has positive probability outside every finite disc. Thus finite agreement of many moments does not imply equality of distributions. As $n$ grows, any fixed moment eventually lies in the exact-agreement range, which is why a limiting Gaussian law remains possible.
\end{solution}

\begin{solution}{ex:v10exp-ch11-weight27}
\idea A determinant twist shifts all weights without changing the dimension, while a central scalar sees their total degree.
\writeup The dimension factors for the pairs $(1,2),(1,3),(2,3)$ are
\[
 3-1+1=3,\qquad\frac{3-(-1)+2}{2}=3,\qquad1-(-1)+1=3.
\]
Thus the dimension is $27$. Adding one to every weight entry gives $(4,2,0)$, so $V_{(3,1,-1)}=\det^{-1}\otimes V_{(4,2,0)}$ up to isomorphism. The determinant inverse has dimension one and does not change $27$. The weight sum is $3$, so $zI$ acts as $z^3I$ and has character $27z^3$. This sum is the central exponent, not a count of basis vectors or a dimension.
\end{solution}

\begin{solution}{ex:v10exp-ch11-matrixcov}
\idea Matrix-coefficient orthogonality computes the covariance even though all entries of a unitary matrix are constrained together.
\writeup Central rotation makes each entry have mean zero, hence makes the trace have mean zero. The defining representation is irreducible, by its identification with $V_{(1,0,\ldots,0)}$. Schur orthogonality gives
$\mathbb E[U_{ji}\overline{U_{lk}}]=\delta_{jl}\delta_{ik}/n$. Expanding both traces and summing yields
\[
 \mathbb E[\operatorname{tr}(AU)\overline{\operatorname{tr}(BU)}]
  =\frac1n\sum_{i,j}A_{ij}\overline{B_{ij}}
  =\frac1n\operatorname{tr}(AB^*).
\]
Taking $B=A$ proves the unbiased-estimate assertion. If $A=E_{ij}$ is a matrix unit, the trace selects $U_{ji}$, whose squared mean is $1/n$; the target norm square is $1$. For $A=I$, the target norm square is $n$ and the trace variance is $1$, agreeing with Proposition~\ref{prop:powertraces}. No independent-entry model was used: the two-entry integral was supplied by representation orthogonality. For $n=1$, the same calculation is ordinary circle integration.
\end{solution}

\begin{solution}{ex:v10exp-ch11-dependent}
\idea Orthogonality removes one mixed covariance; independence would impose many more product identities.
\writeup The power-trace covariance formula gives $\mathbb E[P\overline Q]=0$, $\mathbb E|P|^2=1$ and $\mathbb E|Q|^2=2$. In two variables,
\[
 p_1p_2=(t_1+t_2)(t_1^2+t_2^2)
   =t_1^3+t_1^2t_2+t_1t_2^2+t_2^3=s_{(3,0)}.
\]
This is an irreducible character, so its squared norm is $1$. Thus $\mathbb E|PQ|^2=1$, whereas independence would require $1\cdot2=2$. The variables are not independent.

For three variables, expanding the complete symmetric polynomial $h_3=s_{(3,0,0)}$ shows that $p_1p_2=h_3-e_3$. The monomials of types $t_i^3$ and $t_i^2t_j$ occur once in both $p_1p_2$ and $h_3$, while $h_3$ has the additional monomial $t_1t_2t_3=e_3$. These are two distinct irreducible characters, so the squared norm is $1+1=2$. Agreement with the product of variances at size three is just this one identity. It does not prove independence; independence asks for every product of suitable test functions to factor.
\end{solution}

\begin{solution}{ex:v10exp-ch11-fourthpowertrace}
\idea The signed alternant additions give a finite character expansion; orthogonality squares its coefficients rather than its dimensions.
\writeup Multiplying $a_\delta$ by $p_2$ raises one exponent by two. The first position contributes $(2)$; the second crosses one preceding exponent and contributes $-(1,1)$; the later positions collide. Thus $p_2=s_{(2)}-s_{(1,1)}$, omitting a shape if the dimension is too small.

For a general partition, the same raised-exponent calculation has three outcomes. A new exponent that stays in order adds a horizontal domino, with positive sign. Equality with an old exponent gives zero. Crossing one neighbouring exponent gives a vertical domino, with a negative sign; crossing farther forces a collision. Applied to $(2)$, the allowed terms are $(4),(2,2)$ and $-(2,1,1)$. Applied to $(1,1)$, they are $(3,1),-(2,2)$ and $-(1^4)$. Consequently
\[
\begin{split}
 p_2s_{(2)}&=s_{(4)}+s_{(2,2)}-s_{(2,1,1)},\\
 p_2s_{(1,1)}&=s_{(3,1)}-s_{(2,2)}-s_{(1^4)}.
\end{split}
\]
Subtracting proves the required expansion. In low dimension, discard terms whose row counts exceed $n$; the same alternant calculation on the remaining $n$ exponents gives precisely the surviving terms. At $n=1$ only $s_{(4)}=t_1^4$ remains.

Orthogonality now gives
\[
 \mathbb E|\operatorname{tr}U^2|^4=
 \begin{cases}1,&n=1,\\6,&n=2,\\7,&n=3,\\8,&n\ge4.\end{cases}
\]
For example, at size two the coefficients are $1,-1,2$, so the answer is $1+1+4=6$. A centred standard complex Gaussian $Z$ has $\mathbb E|Z|^4=2$; multiplying it by $\sqrt2$ makes variance $2$ and fourth moment $8$. The exact value stabilises at that number once $n\ge4$, though the finite trace is bounded and is not itself exactly Gaussian. The quoted limit concerns a fixed list of power traces as matrix size grows. No finite list of matching moments alone proves that distributional claim.
\end{solution}

\section*{Sources and further reading}
\addcontentsline{toc}{section}{Sources and further reading}

I.~Schur, ``Neue Anwendungen der Integralrechnung auf Probleme der Invariantentheorie'', \emph{Sitzungsber.\ Preuss.\ Akad.\ Wiss.} (1924), 189--208, 297--321, 346--355; H.~Weyl, ``Theorie der Darstellung kontinuierlicher halb-einfacher Gruppen durch lineare
Transformationen'', \emph{Math.\ Z.} 23 (1925), 271--309, and 24 (1926), 328--395, 789--791; F.~Peter and H.~Weyl, ``Die Vollständigkeit der primitiven Darstellungen einer geschlossenen kontinuierlichen Gruppe'', \emph{Math.\ Ann.} 97 (1927), 737--755;
H.~Weyl, ``Invariants'', \emph{Duke Math.\ J.} 5 (1939), 489--502, for the quotation; H.~Weyl, \emph{The Classical Groups} (Princeton University Press, 1939). P.~Diaconis and M.~Shahshahani, ``On the eigenvalues of random matrices'', \emph{J.\ Appl.\ Probab.} 31A (1994), 49--62;
E.~M.~Rains, ``Increasing subsequences and the classical groups'', \emph{Electron.\ J.\ Combin.} 5 (1998), R12; I.~M.~Gessel, ``Symmetric functions and P-recursiveness'', \emph{J.\ Combin.\ Theory Ser.\ A} 53 (1990), 257--285; J.~P.~Keating and N.~C.~Snaith, ``Random matrix
theory and $\zeta(1/2+it)$'', \emph{Comm.\ Math.\ Phys.} 214 (2000), 57--89; M.~Gell-Mann, \emph{The Eightfold Way: a Theory of Strong Interaction Symmetry}, Caltech report CTSL-20 (1961).

For more: T.~Bröcker and T.~tom Dieck, \emph{Representations of Compact Lie Groups} (Springer, 1985), Chapters~II--VI; D.~Bump, \emph{Lie Groups} (Springer, 2nd edition 2013), for Schur--Weyl duality and the Keating--Snaith formula; I.~G.~Macdonald, \emph{Symmetric Functions and
Hall Polynomials} (Oxford, 2nd edition 1995), Chapter~I; E.~S.~Meckes, \emph{The Random Matrix Theory of the Classical Compact Groups} (Cambridge University Press, 2019).
